\documentclass{article}
\usepackage{amsmath,amssymb,amsthm}
\usepackage{algorithm,algorithmic}
\usepackage{enumitem}
\usepackage{hyperref}
\usepackage{geometry}
\geometry{margin=1in}
\newtheorem{theorem}{Theorem}
\newtheorem{lemma}{Lemma}
\newtheorem{assumption}{Assumption}
\newcommand{\E}{\mathbb{E}}
\newcommand{\R}{\mathbb{R}}

\begin{document}

\title{Quantized Stochastic Gradient Descent (QSGD) \\ Detailed Algorithmic Description and Convergence Proof}
\author{}
\date{}
\maketitle

%%%%%%%%%%%%%%%%%%%%%%%%%%%%%%%%%%%%%%%%%%%%%%%%%%%%%%%%%%%%%%%%%%%%%%%%%%%%%%%
\section*{Algorithm Name}
\textbf{Quantized Stochastic Gradient Descent (QSGD)}

%%%%%%%%%%%%%%%%%%%%%%%%%%%%%%%%%%%%%%%%%%%%%%%%%%%%%%%%%%%%%%%%%%%%%%%%%%%%%%%
\section*{Mathematical Formulation}
Let $f:\R^{d}\rightarrow \R$ be a (possibly non‑convex) smooth objective.  
At iteration $t$ we draw a stochastic sample $\xi_t$ and compute the stochastic
gradient
\[
g_t \;:=\; \nabla f(w_t;\xi_t)\in\R^{d}.
\]
A (random) quantization operator $Q:\R^d\to\R^d$ satisfies

\[
\boxed{\;\E[Q(g)\mid g]=g\;}\qquad\text{(unbiasedness)} \tag{1}
\]
and
\[
\boxed{\;\E\!\big[\|Q(g)\|^{2}\mid g\big]\le C\|g\|^{2}\;}\qquad\text{(second‑moment bound)} \tag{2}
\]
for a constant $C\ge 1$ (the \emph{quantization quality factor}).

The update rule of QSGD is
\[
\boxed{\; w_{t+1}=w_t-\eta_t\,Q(g_t)\;}, \tag{3}
\]
where $\{\eta_t\}_{t\ge0}$ is a stepsize (learning‑rate) schedule.

%%%%%%%%%%%%%%%%%%%%%%%%%%%%%%%%%%%%%%%%%%%%%%%%%%%%%%%%%%%%%%%%%%%%%%%%%%%%%%%
\section*{Pseudocode}
\begin{algorithm}[H]
\caption{Quantized Stochastic Gradient Descent (QSGD)}
\begin{algorithmic}[1]
\STATE \textbf{Input:} initial point $w_0\in\R^{d}$, stepsize schedule $\{\eta_t\}$, 
quantizer $Q(\cdot)$ satisfying (1)–(2), maximum iterations $T$.
\FOR{$t=0,1,\dots,T-1$}
    \STATE Sample a minibatch (or a random data point) $\xi_t$.
    \STATE Compute stochastic gradient $g_t=\nabla f(w_t;\xi_t)$.
    \STATE Quantize: $\tilde g_t = Q(g_t)$.
    \STATE Update: $w_{t+1}=w_t-\eta_t \tilde g_t$.
\ENDFOR
\STATE \textbf{Return} $w_T$ (or optionally the averaged iterate $\bar w_T=\frac1T\sum_{t=0}^{T-1}w_t$).
\end{algorithmic}
\end{algorithm}

%%%%%%%%%%%%%%%%%%%%%%%%%%%%%%%%%%%%%%%%%%%%%%%%%%%%%%%%%%%%%%%%%%%%%%%%%%%%%%%
\section*{Dry Run Example}
Consider the scalar quadratic $f(w)=(w-3)^2$, thus $\nabla f(w)=2(w-3)$.  
Choose a deterministic ``quantizer'' that simply adds zero‑mean noise:
\[
Q(g)=g+\zeta,\qquad \E[\zeta]=0,\;\E[\zeta^2]=\sigma_Q^2,\;\sigma_Q^2=0.01 .
\]
Hence (1)–(2) hold with $C=1+\sigma_Q^2/ g^2$; for this simple example we treat $C\approx1$.

Parameters: $w_0=0$, constant stepsize $\eta=0.1$, $T=3$.

\begin{center}
\begin{tabular}{c|c|c|c}
$t$ & $g_t=2(w_t-3)$ & $Q(g_t)=g_t+\zeta_t$ & $w_{t+1}=w_t-\eta Q(g_t)$ \\ \hline
0 & $2(0-3)=-6$ & $-6+ \zeta_0$ & $w_1=0-0.1(-6+\zeta_0)=0.6-0.1\zeta_0$ \\
1 & $2(w_1-3)=2(0.6-0.1\zeta_0-3)=-4.8-0.2\zeta_0$ &
$-4.8-0.2\zeta_0+\zeta_1$ &
$w_2=w_1-0.1(-4.8-0.2\zeta_0+\zeta_1)$ \\
2 & $\displaystyle 2(w_2-3)$ &
$Q(g_2)=g_2+\zeta_2$ &
$w_3=w_2-0.1\,Q(g_2)$
\end{tabular}
\end{center}
Taking expectations (recall $\E[\zeta_i]=0$) gives the deterministic trajectory
\[
\begin{aligned}
\E[w_1] &= 0.6, &
\E[w_2] &= 0.6 +0.1\cdot4.8 = 1.08, &
\E[w_3] &= 1.08 +0.1\cdot3.84 = 1.464 .
\end{aligned}
\]
Corresponding objective values (in expectation) are
\[
\E[f(w_1)]= (0.6-3)^2 =5.76,\;
\E[f(w_2)]=(1.08-3)^2=3.6864,\;
\E[f(w_3)]=(1.464-3)^2=2.366.
\]
Thus each iteration reduces the expected loss, illustrating the effect of the unbiased quantizer.

%%%%%%%%%%%%%%%%%%%%%%%%%%%%%%%%%%%%%%%%%%%%%%%%%%%%%%%%%%%%%%%%%%%%%%%%%%%%%%%
\section*{Assumptions}
\begin{assumption}[Smoothness]
$f$ is $L$‑smooth, i.e.
\[
\forall x,y\in\R^{d}:\quad
f(y)\le f(x)+\langle\nabla f(x),y-x\rangle+\frac{L}{2}\|y-x\|^{2}. \tag{A1}
\]
\end{assumption}
\begin{assumption}[Unbiased Stochastic Gradient and Bounded Variance]
For every $t$,
\[
\E\!\big[g_t\mid w_t\big]=\nabla f(w_t),\qquad
\E\!\big[\|g_t-\nabla f(w_t)\|^{2}\mid w_t\big]\le\sigma^{2}. \tag{A2}
\]
\end{assumption}
\begin{assumption}[Quantizer Properties]
The quantizer $Q$ satisfies (1)–(2) with a constant $C\ge1$ independent of $t$ and $w_t$.
\tag{A3}
\end{assumption}
\begin{assumption}[Stepsize]
We use a constant stepsize 
\[
\eta_t\equiv \eta\le\frac{1}{L\sqrt{C}}. \tag{A4}
\]
\end{assumption}

%%%%%%%%%%%%%%%%%%%%%%%%%%%%%%%%%%%%%%%%%%%%%%%%%%%%%%%%%%%%%%%%%%%%%%%%%%%%%%%
\section*{Main Convergence Theorem}
\begin{theorem}[Expected Gradient Norm Convergence]\label{thm:main}
Under Assumptions (A1)–(A4) and after $T$ iterations of QSGD, the
averaged iterate $\bar w_T:=\frac1T\sum_{t=0}^{T-1}w_t$ satisfies
\[
\frac{1}{T}\sum_{t=0}^{T-1}\E\!\big[\|\nabla f(w_t)\|^{2}\big]
\;\le\;
\frac{2\big(f(w_0)-f^{\star}\big)}{\eta T}
\;+\;
\eta L C\,\sigma^{2},                                            \tag{4}
\]
where $f^{\star}=\inf_{w}f(w)$.
Consequently, choosing $\eta = \Theta\!\big(\tfrac{1}{\sqrt{T}}\big)$ yields
\[
\frac{1}{T}\sum_{t=0}^{T-1}\E\!\big[\|\nabla f(w_t)\|^{2}\big]
= O\!\left(\frac{1}{\sqrt{T}}\right).
\]
\end{theorem}

\begin{theorem}[Special Case: Exact Gradients]
If $g_t=\nabla f(w_t)$ (i.e., $\sigma=0$) then (4) reduces to
\[
\frac{1}{T}\sum_{t=0}^{T-1}\E\!\big[\|\nabla f(w_t)\|^{2}\big]
\le \frac{2\big(f(w_0)-f^{\star}\big)}{\eta T},
\]
so the method attains the deterministic $O(1/T)$ rate.
\end{theorem}

%%%%%%%%%%%%%%%%%%%%%%%%%%%%%%%%%%%%%%%%%%%%%%%%%%%%%%%%%%%%%%%%%%%%%%%%%%%%%%%
\section*{Theoretical Properties}
\begin{itemize}[leftmargin=*]
\item \textbf{Stability.}  The bound (4) holds for any $C\ge1$; a finer quantizer (smaller $C$) directly reduces the variance term $\eta L C\sigma^{2}$.
\item \textbf{Hyperparameter Sensitivity.}  The admissible stepsize is limited by $\eta\le 1/(L\sqrt{C})$; as $C$ grows, the safe stepsize shrinks proportionally to $1/\sqrt{C}$.
\item \textbf{Comparison to Plain SGD.}  Setting $C=1$ recovers the standard SGD guarantee.  Quantization therefore incurs at most a factor $C$ slowdown in the variance term.
\item \textbf{Non‑convexity.}  No convexity is required; the guarantee is on the \emph{average} squared gradient norm, the standard stationarity measure for non‑convex optimization.
\end{itemize}

%%%%%%%%%%%%%%%%%%%%%%%%%%%%%%%%%%%%%%%%%%%%%%%%%%%%%%%%%%%%%%%%%%%%%%%%%%%%%%%
\section*{Preliminaries (Definitions and Lemmas)}
\begin{lemma}[Smoothness Inequality]\label{lem:smooth}
If $f$ is $L$‑smooth then for any $x,y\in\R^{d}$,
\[
f(y)\le f(x)+\langle\nabla f(x),y-x\rangle+\frac{L}{2}\|y-x\|^{2}.
\]
\end{lemma}
\begin{lemma}[Unbiasedness of Quantizer]\label{lem:unbiased}
For any random vector $g$, $\E[Q(g)\mid g]=g$.
\end{lemma}
\begin{lemma}[Second‑Moment Bound]\label{lem:second}
For any $g$, $\E[\|Q(g)\|^{2}\mid g]\le C\|g\|^{2}$.
\end{lemma}

%%%%%%%%%%%%%%%%%%%%%%%%%%%%%%%%%%%%%%%%%%%%%%%%%%%%%%%%%%%%%%%%%%%%%%%%%%%%%%%
\section*{Main Proof (Step‑by‑Step Derivation)}
We prove Theorem~\ref{thm:main}.  Throughout we condition on the filtration
$\mathcal{F}_t:=\sigma(w_0,\xi_0,\dots,\xi_{t-1})$ generated by all randomness up to
iteration $t$.

\paragraph{Step 1 (Smoothness expansion).}
By Lemma~\ref{lem:smooth} with $x=w_t$, $y=w_{t+1}=w_t-\eta Q(g_t)$,
\[
\begin{aligned}
f(w_{t+1})
&\le f(w_t)+\langle\nabla f(w_t),\,w_{t+1}-w_t\rangle
   +\frac{L}{2}\|w_{t+1}-w_t\|^{2}\\
&= f(w_t)-\eta\langle\nabla f(w_t),\,Q(g_t)\rangle
   +\frac{L\eta^{2}}{2}\|Q(g_t)\|^{2}.          \tag{5}
\end{aligned}
\]
\paragraph{Step 2 (Take conditional expectation).}
Apply $\E[\cdot\mid\mathcal{F}_t]$ to (5):
\[
\E\!\big[f(w_{t+1})\mid\mathcal{F}_t\big]
\le f(w_t)-\eta\,\E\!\big[\langle\nabla f(w_t),Q(g_t)\rangle\mid\mathcal{F}_t\big]
   +\frac{L\eta^{2}}{2}\E\!\big[\|Q(g_t)\|^{2}\mid\mathcal{F}_t\big]. \tag{6}
\]
\paragraph{Step 3 (Use unbiasedness of $Q$).}
By Lemma~\ref{lem:unbiased} and Assumption (A2),
\[
\E\!\big[\langle\nabla f(w_t),Q(g_t)\rangle\mid\mathcal{F}_t\big]
= \big\langle\nabla f(w_t),\E[Q(g_t)\mid\mathcal{F}_t]\big\rangle
= \big\langle\nabla f(w_t),\E[g_t\mid\mathcal{F}_t]\big\rangle
= \|\nabla f(w_t)\|^{2}.                                            \tag{7}
\]
\paragraph{Step 4 (Bound the second moment of the quantized gradient).}
Using Lemma~\ref{lem:second} and the variance decomposition,
\[
\begin{aligned}
\E\!\big[\|Q(g_t)\|^{2}\mid\mathcal{F}_t\big]
&\le C\,\E\!\big[\|g_t\|^{2}\mid\mathcal{F}_t\big] \\
&= C\,\big\|\E[g_t\mid\mathcal{F}_t]\big\|^{2}
   + C\,\E\!\big[\|g_t-\E[g_t\mid\mathcal{F}_t]\|^{2}\mid\mathcal{F}_t\big]  \\
&= C\,\|\nabla f(w_t)\|^{2}+C\,\E\!\big[\|g_t-\nabla f(w_t)\|^{2}\mid\mathcal{F}_t\big] \\
&\le C\,\|\nabla f(w_t)\|^{2}+C\sigma^{2}.                         \tag{8}
\end{aligned}
\]
\paragraph{Step 5 (Plug (7) and (8) into (6)).}
\[
\E\!\big[f(w_{t+1})\mid\mathcal{F}_t\big]
\le f(w_t)-\eta\|\nabla f(w_t)\|^{2}
   +\frac{L\eta^{2}C}{2}\|\nabla f(w_t)\|^{2}
   +\frac{L\eta^{2}C}{2}\sigma^{2}.                               \tag{9}
\]
\paragraph{Step 6 (Rearrange).}
Collect the terms containing $\|\nabla f(w_t)\|^{2}$:
\[
\E\!\big[f(w_{t+1})\mid\mathcal{F}_t\big]
\le f(w_t)-\underbrace{\Big(\eta-\frac{L\eta^{2}C}{2}\Big)}_{\displaystyle\alpha}
   \|\nabla f(w_t)\|^{2}
   +\frac{L\eta^{2}C}{2}\sigma^{2}.                               \tag{10}
\]
\paragraph{Step 7 (Choose stepsize to make $\alpha>0$).}
By Assumption (A4), $\eta\le 1/(L\sqrt{C})$, hence
\[
\eta-\frac{L\eta^{2}C}{2}\;\ge\;\frac{\eta}{2}\;=: \alpha>0.          \tag{11}
\]
Thus (10) simplifies to
\[
\E\!\big[f(w_{t+1})\mid\mathcal{F}_t\big]
\le f(w_t)-\frac{\eta}{2}\|\nabla f(w_t)\|^{2}
   +\frac{L\eta^{2}C}{2}\sigma^{2}.                               \tag{12}
\]

\paragraph{Step 8 (Take full expectation and sum).}
Taking expectation over $\mathcal{F}_t$ and summing $t=0,\dots,T-1$ yields
\[
\begin{aligned}
\E[f(w_T)]
&\le f(w_0)-\frac{\eta}{2}\sum_{t=0}^{T-1}\E\!\big[\|\nabla f(w_t)\|^{2}\big]
   +\frac{L\eta^{2}C}{2}\sigma^{2}T.                               \tag{13}
\end{aligned}
\]
Rearranging,
\[
\frac{1}{T}\sum_{t=0}^{T-1}\E\!\big[\|\nabla f(w_t)\|^{2}\big]
\le \frac{2\big(f(w_0)-\E[f(w_T)]\big)}{\eta T}
   +L\eta C\sigma^{2}.                                            \tag{14}
\]
Since $f(w_T)\ge f^{\star}$,
\[
\frac{1}{T}\sum_{t=0}^{T-1}\E\!\big[\|\nabla f(w_t)\|^{2}\big]
\le \frac{2\big(f(w_0)-f^{\star}\big)}{\eta T}
   +L\eta C\sigma^{2},                                            \tag{15}
\]
which is exactly the bound (4) claimed in Theorem~\ref{thm:main}.  ∎

\paragraph{Step 9 (Choosing $\eta$ for the optimal rate).}
Set $\eta = \sqrt{\frac{2(f(w_0)-f^{\star})}{L C\sigma^{2}T}}$,
which respects $\eta\le 1/(L\sqrt{C})$ for all $T\ge1$.  Substituting into (15)
gives
\[
\frac{1}{T}\sum_{t=0}^{T-1}\E\!\big[\|\nabla f(w_t)\|^{2}\big]
= O\!\left(\frac{1}{\sqrt{T}}\right).
\]

%%%%%%%%%%%%%%%%%%%%%%%%%%%%%%%%%%%%%%%%%%%%%%%%%%%%%%%%%%%%%%%%%%%%%%%%%%%%%%%
\section*{Rate Derivation (With Constants)}
From (15) we can write the explicit constant‑aware bound:
\[
\boxed{
\frac{1}{T}\sum_{t=0}^{T-1}\E\!\big[\|\nabla f(w_t)\|^{2}\big]
\le
\underbrace{\frac{2\big(f(w_0)-f^{\star}\big)}{\eta T}}_{\text{optimization term}}
\;+\;
\underbrace{L C\eta\sigma^{2}}_{\text{variance term}} }.
\]
If $\eta = \Theta(1/\sqrt{T})$ both terms scale as $O(1/\sqrt{T})$.

%%%%%%%%%%%%%%%%%%%%%%%%%%%%%%%%%%%%%%%%%%%%%%%%%%%%%%%%%%%%%%%%%%%%%%%%%%%%%%%
\section*{Special Cases}
\begin{enumerate}[leftmargin=*]
\item \textbf{Exact Gradients ($\sigma=0$).}  The variance term vanishes, and (15) reduces to
\[
\frac{1}{T}\sum_{t=0}^{T-1}\E\!\big[\|\nabla f(w_t)\|^{2}\big]
\le\frac{2\big(f(w_0)-f^{\star}\big)}{\eta T},
\]
recovering the deterministic $O(1/T)$ rate.
\item \textbf{Perfect Quantization ($C=1$).}  The bound matches the classic SGD guarantee.
\item \textbf{Coarse Quantization ($C\gg1$).}  The stepsize must be reduced according
to (A4), and the variance term is inflated by $C$, reflecting the extra noise introduced by quantization.
\end{enumerate}

%%%%%%%%%%%%%%%%%%%%%%%%%%%%%%%%%%%%%%%%%%%%%%%%%%%%%%%%%%%%%%%%%%%%%%%%%%%%%%%
\section*{Discussion}
The proof shows that unbiased quantization only inflates the stochastic variance by the factor $C$ and forces a more conservative stepsize.  Consequently, QSGD retains the $O(1/\sqrt{T})$ stationarity rate of vanilla SGD for non‑convex smooth objectives, while dramatically reducing communication cost.  The analysis is tight in the sense that both the optimization term and the variance term appear in the final bound and cannot be simultaneously improved without stronger assumptions (e.g., bounded gradients or variance reduction).

\end{document}